\section{Problemstellung}

Im Labor für Maschinenelemente sind mechanische Verbindungs- und Maschinenelemente ausgestellt. Einige dieser Bauteile sind in einer Glasvitrine ausgestellt (siehe Abb. \ref{fig:schraubenGesamt}). Das Problem hierbei ist, einige der mechanischen Verbindungselemente haben ein hohes Gewicht und könnten beim unbeabsichtigt zu starkem Ablegen das Glas beschädigen. Zudem liegen die Bauteile ungeordnet frei herum. 

\begin{figure}[H]
	\centering
	\includegraphics[width=0.7\linewidth]{schraubenGesamt.png}
	\caption{Verbindungselemente liegen auf der Glasplatte}
	\label{fig:schraubenGesamt}
\end{figure}

Zusätzlich befinden sich noch axiale Wellensicherungselemente in der Glasvitrine, siehe dazu Abb. \ref{fig:sicherungsringe}. \\
In Abb. \ref{fig:sicherungsringe} liegen die Sicherungselemente auf einer Holzplatte für eine bessere Sichtbarkeit. Diese liegen ebenfalls regellos in der Vitrine und sind schlecht sichtbar. Zudem lässt sich die spätere Funktion als  Maschinenelement auf den ersten Blick nicht erkennen.

\begin{figure}[H]
	\centering
	\includegraphics[width=0.7\linewidth]{wellensicherungen.png}
	\caption{Axiale Wellensicherungen sind ungeordnet verteilt}
	\label{fig:sicherungsringe}
\end{figure}